\subsection{Graphing a Quadratic Function in Vertex Form} 
\begin{fact}
If a quadratic function is written in vertex form as $$f(x)=a(x-h)^2+k$$
\begin{itemize*}
\item Its vertex is at $(h,k)$.
\item If $a>0$, the parabola opens upward. If $a<0$, the parabola opens downward.
\item For any number $d$, the points where $x=h+d$ and $x=h-d$ have the same $y$-coordinate (because of \makeblank[1.5in]).
\end{itemize*}
\end{fact}
To graph a parabola, we need the vertex and two points to either side. When the function is given in vertex form, this is easy to find.
\begin{example}
Sketch the graph of the function $h(x)=2(x+3)^2+1$.
\begin{lpic}{}{13}
\setgraph{6}{2}{1}{10}
\GraphInSet
\regtextleft{0.25}{-1}{Vertex:};
\regtext{1}{-4}{$x$};
\regtext{4}{-4}{$y=h(x)$};
\draw (0,-4) -- (6,-4);
\draw (2,-3) -- (2,-9);
\end{lpic}
\end{example}